\documentclass[11pt]{article}
\usepackage[T1]{fontenc}
\usepackage[margin=1in]{geometry}
\usepackage{amsmath,amssymb,amsthm,mathtools}
\usepackage{graphicx}
\usepackage{microtype}
\usepackage{xcolor}
\usepackage[colorlinks=true,linkcolor=blue!55!black,citecolor=blue!55!black,urlcolor=blue!65!black]{hyperref}

\newtheorem{theorem}{Theorem}
\newtheorem{lemma}[theorem]{Lemma}
\newcommand{\A}{\mathcal A}
\newcommand{\B}{\mathcal B}
\newcommand{\D}{\mathbb D}
\newcommand{\C}{\mathbb C}
\emergencystretch=3em

\title{Square-zero local Banach algebras satisfying von Neumann's inequality\\
are isometrically operator algebras}
\author{Partial-result packet for arXiv:1804.04062}
\date{27 August 2026}

\begin{document}
\maketitle

\begin{center}
\fbox{\begin{minipage}{0.91\textwidth}
\textbf{Status: candidate partial result, likely valid.}
The general question remains open.  The result below gives a complete
affirmative answer for every normalized unital algebra
$\A=\C1\oplus X$ whose closed codimension-one ideal satisfies $X^2=0$,
without a dimension assumption.
\end{minipage}}
\end{center}

\section{Source question}

Badea, Crouzeix, and Klaja ask whether every normalized unital complex Banach
algebra satisfying the one-variable von Neumann inequality is isometrically
isomorphic to a closed algebra of Hilbert-space operators
\cite[p.~4]{BCK}.  Explicitly, their hypothesis is
\[
  \|a\|\leq 1 \quad\Longrightarrow\quad
  \|p(a)\|\leq \max_{|z|\leq1}|p(z)|
  \qquad(p\in\C[z]).
\]
The source passage is reproduced below.

\begin{center}
\includegraphics[width=\textwidth]{figures/open_problem_crop.png}
\end{center}

The non-unital analogue is false, but the normalized unit is the obstruction
to simply adjoining an identity to those counterexamples.  The source's
general unital question was still recorded as open in 2025
\cite[p.~31]{CT}.

\section{Partial theorem}

\subsection*{Proof intuition}

Write an element as $\lambda1+x$, with $x^2=0$.  Holomorphic functional
calculus then remembers only the value and first derivative of a disk map:
$f(\lambda1+x)=f(\lambda)1+f'(\lambda)x$.  A disk automorphism sending
$\lambda$ to zero expands the radical direction by exactly
$(1-|\lambda|^2)^{-1}$.  Its inverse gives the reverse inclusion, so the entire
unit ball is forced by its central slice.  That forced ball is precisely the
contraction ball of a $2\times2$ upper-triangular block operator.  Embedding the
Banach space $X$ isometrically as diagonal Hilbert-space operators completes
the representation.

\begin{theorem}[square-zero local class]\label{thm:main}
Let $\A$ be a complex unital Banach algebra with $\|1\|=1$.  Suppose
\[
  \A=\C1\oplus X,
\]
where $X$ is a closed two-sided ideal and $X^2=\{0\}$.  If $\A$ satisfies von
Neumann's inequality, then, writing $q(x)=\|x\|_{\A}$ for $x\in X$, one has
\begin{equation}\label{eq:ball}
  \|\lambda1+x\|_{\A}\leq1
  \quad\Longleftrightarrow\quad
  |\lambda|^2+q(x)\leq1.
\end{equation}
Consequently,
\begin{equation}\label{eq:norm}
  \|\lambda1+x\|_{\A}
  =\frac{q(x)+\sqrt{q(x)^2+4|\lambda|^2}}{2}.
\end{equation}
Moreover, there are a Hilbert space $H$ and a linear isometry
$J:X\to\B(H)$ such that
\begin{equation}\label{eq:rep}
  \Pi(\lambda1+x)=
  \begin{pmatrix}
    \lambda I_H&Jx\\
    0&\lambda I_H
  \end{pmatrix}
\end{equation}
is a unital isometric algebra homomorphism
$\A\to\B(H\oplus H)$.  Its range is closed, so $\A$ is an operator algebra.
\end{theorem}

\section{Proof}

\begin{lemma}[disk-algebra functional calculus]\label{lem:disk}
If $a$ is in the closed unit ball of a unital Banach algebra satisfying von
Neumann's inequality, then
\[
  \|f(a)\|\leq\|f\|_{\infty,\overline\D}
  \qquad(f\in A(\D)).
\]
In particular this holds for every automorphism of $\D$.
\end{lemma}

\begin{proof}
Choose polynomials converging uniformly to $f$ on $\overline\D$.  Applying the
given inequality to differences shows that their values at $a$ are Cauchy and
that the limit has norm at most $\|f\|_\infty$.  For a function holomorphic on
a neighborhood of $\overline\D$, this limit agrees with the usual holomorphic
functional calculus.
\end{proof}

\begin{lemma}[forced fibers of the unit ball]\label{lem:fibers}
Under the assumptions of Theorem~\ref{thm:main}, formula~\eqref{eq:ball}
holds.
\end{lemma}

\begin{proof}
Let $a=\lambda1+x$ be in the closed unit ball.  The quotient character
$\A\to\A/X\cong\C$ has norm one, so $|\lambda|\leq1$.

First suppose $|\lambda|<1$ and set
\[
  \phi_\lambda(z)=\frac{z-\lambda}{1-\overline\lambda z}.
\]
Because $X^2=0$, polynomial calculus, followed by uniform approximation,
gives
\[
 f(\lambda1+x)=f(\lambda)1+f'(\lambda)x
\]
for every function used here.  Lemma~\ref{lem:disk} therefore yields
\[
 \left\|\phi_\lambda(a)\right\|
 =\left\|\frac{x}{1-|\lambda|^2}\right\|\leq1,
\]
and hence $q(x)\leq1-|\lambda|^2$.

Conversely, assume $q(x)\leq1-|\lambda|^2$.  Then
$y=x/(1-|\lambda|^2)$ is in the unit ball.  The inverse disk automorphism
\[
  \psi_\lambda(z)=\frac{z+\lambda}{1+\overline\lambda z}
\]
satisfies $\psi_\lambda(0)=\lambda$ and
$\psi_\lambda'(0)=1-|\lambda|^2$.  Thus
\[
  \psi_\lambda(y)=\lambda1+x
\]
is again in the unit ball by Lemma~\ref{lem:disk}.

It remains to treat $|\lambda|=1$.  For each $n\geq1$, von Neumann's
inequality applied to $z^n$ gives
\[
  \left\|\lambda^{-n}a^n\right\|
  =\left\|1+n\lambda^{-1}x\right\|\leq1.
\]
Consequently $nq(x)\leq2$, for every $n$, and $x=0$.  This proves both
directions of~\eqref{eq:ball}.
\end{proof}

The norm formula~\eqref{eq:norm} now follows by taking the Minkowski
functional of~\eqref{eq:ball}: the least $t>0$ for which
$|\lambda|^2/t^2+q(x)/t\leq1$ is the positive root of
$t^2-q(x)t-|\lambda|^2=0$.

\begin{lemma}[upper-triangular contraction test]\label{lem:block}
For a Hilbert-space operator $S$ and $\lambda\in\C$, put
\[
 M_{\lambda,S}=\begin{pmatrix}\lambda I&S\\0&\lambda I\end{pmatrix}.
\]
Then
\[
 \|M_{\lambda,S}\|\leq1
 \quad\Longleftrightarrow\quad
 |\lambda|\leq1\ \text{ and }\ \|S\|\leq1-|\lambda|^2.
\]
\end{lemma}

\begin{proof}
For $|\lambda|<1$, let $c=1-|\lambda|^2$.  The upper-left corner of
$I-M_{\lambda,S}^*M_{\lambda,S}$ is $cI>0$, and its Schur complement is
\[
 cI-S^*S-\frac{|\lambda|^2}{c}S^*S
 =cI-\frac1cS^*S.
\]
Thus positivity is equivalent to $S^*S\leq c^2I$.  If $|\lambda|=1$,
positivity forces $S=0$; if $|\lambda|>1$, the spectral radius already exceeds
one.
\end{proof}

\begin{proof}[Proof of Theorem~\ref{thm:main}]
Only the representation remains.  Let $K$ be the closed unit ball of $X^*$
and put $H=\ell_2(K)$.  The canonical diagonal map
\[
  (Jx)e_\varphi=\varphi(x)e_\varphi\qquad(\varphi\in K)
\]
is a linear isometry $X\to\B(H)$ by Hahn--Banach.  Since $X^2=0$, direct
block multiplication shows that~\eqref{eq:rep} is a unital algebra
homomorphism.  Lemma~\ref{lem:block}, $\|Jx\|=q(x)$, and
Lemma~\ref{lem:fibers} show that $\Pi$ preserves the closed unit ball exactly.
Therefore $\Pi$ is an isometry.  Its range is closed because $\A$ is complete.
\end{proof}

\section{Verification, scope, and novelty audit}

The block calculation was also checked in 2,000 seeded random complex matrix
cases by \nolinkurl{code/verify_block_contraction.py}; the largest discrepancy
between the computed singular-value norm and formula~\eqref{eq:norm} was at
floating-point roundoff.  This computation is only a regression check; the
Schur-complement argument is the proof.

The theorem covers arbitrary-dimensional (including nonreflexive) square-zero
radicals, but it requires a one-dimensional semisimple quotient and does not
settle the source question for general radicals or multiple characters.

For novelty, the run searched its solution, attempt, proof-gap, and registry
indexes for arXiv:1804.04062 and the core terms.  Bounded exact-phrase searches
for the unital question, two-dimensional algebras, dual numbers, square-zero
ideals, and radical-square-zero algebras recovered Nakazi--Yamamoto's complete
two-dimensional theorem \cite{NY} and Paulsen--Solazzo's semisimple
$\C^k$ analysis \cite{PS}, but no arbitrary-dimensional theorem with
$\A=\C1\oplus X$, $X^2=0$.  Chirvasitu--Thompson still describe the general
unital characterization as unknown in 2025 \cite{CT}.  Novelty is therefore
plausible but provisional; a specialist literature review is recommended.

\paragraph{Human-review focus.}
Verify the passage from polynomial to disk-algebra functional calculus, the
two inclusions in the fiber formula, and whether the arbitrary-dimensional
square-zero statement appears under different terminology in the operator
algebra literature.

\begin{thebibliography}{9}
\bibitem{BCK}
C.~Badea, M.~Crouzeix, and H.~Klaja,
\emph{Spectral sets and operator radii}, arXiv:1804.04062 (2018).

\bibitem{NY}
T.~Nakazi and T.~Yamamoto,
\emph{Two dimensional commutative Banach algebras and von Neumann inequality},
Linear Algebra Appl. 341 (2002), 273--290.

\bibitem{PS}
V.~I.~Paulsen and J.~P.~Solazzo,
\emph{Interpolation and balls in $\C^k$},
J. Operator Theory 60 (2008), 379--398.

\bibitem{CT}
A.~Chirvasitu and I.~Thompson,
\emph{Local presentability and monadicity of forgetful functors for operator
algebraic categories}, arXiv:2507.23152 (2025).
\end{thebibliography}

\end{document}
